% !TeX program = xelatex
\documentclass[11pt,a4paper,openany]{report}

\usepackage[a4paper,margin=2.2cm]{geometry}
\usepackage{amsmath,amssymb,bm}
\usepackage{booktabs}
\usepackage{longtable}
\usepackage{array}
\usepackage{enumitem}
\usepackage{xcolor}
\usepackage{fontspec}
\usepackage{hyperref}
\usepackage{graphicx}
\usepackage{listings}
\usepackage[most]{tcolorbox}
\usepackage{tikz}
\usetikzlibrary{arrows.meta,positioning,fit,calc}

\definecolor{FuseBlue}{HTML}{2563EB}
\definecolor{FuseGreen}{HTML}{16A34A}
\definecolor{FuseOrange}{HTML}{EA580C}
\definecolor{FuseRed}{HTML}{DC2626}
\definecolor{CodeBg}{HTML}{F6F8FA}
\definecolor{FrameGray}{HTML}{D0D7DE}

\setmainfont{Noto Serif CJK SC}
\setsansfont{Noto Sans CJK SC}
\setmonofont{Noto Sans Mono CJK SC}
\XeTeXlinebreaklocale "zh"
\XeTeXlinebreakskip = 0pt plus 1pt
\emergencystretch=2em
\renewcommand{\contentsname}{目录}
\renewcommand{\figurename}{图}
\renewcommand{\tablename}{表}

\hypersetup{
  colorlinks=true,
  linkcolor=FuseBlue,
  citecolor=FuseBlue,
  urlcolor=FuseBlue,
  pdftitle={多传感器融合算法教程：面向 Unitree Go2 RSSI 路径进度项目},
  pdfauthor={Go2 RSSI Project}
}

\setlist[itemize]{leftmargin=2em,itemsep=0.2em,topsep=0.4em}
\setlist[enumerate]{leftmargin=2em,itemsep=0.2em,topsep=0.4em}
\renewcommand{\arraystretch}{1.18}
\setlength{\tabcolsep}{4pt}
\newcolumntype{L}[1]{>{\raggedright\arraybackslash}p{#1}}
\urlstyle{same}

\lstdefinestyle{codestyle}{
  backgroundcolor=\color{CodeBg},
  basicstyle=\ttfamily\small,
  breaklines=true,
  breakatwhitespace=false,
  columns=fullflexible,
  keepspaces=true,
  frame=single,
  rulecolor=\color{FrameGray},
  xleftmargin=0.4em,
  xrightmargin=0.4em,
  aboveskip=0.8em,
  belowskip=0.8em,
  showstringspaces=false,
  tabsize=2
}
\lstdefinelanguage{yaml}{
  keywords={true,false,null,y,n},
  keywordstyle=\color{FuseBlue}\bfseries,
  basicstyle=\ttfamily\small,
  sensitive=false,
  comment=[l]{\#},
  morestring=[b]',
  morestring=[b]"
}
\lstdefinelanguage{bash}{
  keywords={cd,source,export,if,then,else,elif,fi,for,do,done,case,esac,function,exec,true,false},
  keywordstyle=\color{FuseBlue}\bfseries,
  basicstyle=\ttfamily\small,
  sensitive=true,
  comment=[l]{\#},
  morestring=[b]',
  morestring=[b]"
}
\lstset{style=codestyle}

\newtcolorbox{notebox}[1]{
  enhanced,
  breakable,
  colback=blue!3,
  colframe=FuseBlue!65,
  title=#1,
  fonttitle=\bfseries,
  boxrule=0.6pt,
  arc=2mm,
  left=1mm,
  right=1mm,
  top=1mm,
  bottom=1mm
}

\newtcolorbox{warnbox}[1]{
  enhanced,
  breakable,
  colback=orange!6,
  colframe=FuseOrange!75,
  title=#1,
  fonttitle=\bfseries,
  boxrule=0.6pt,
  arc=2mm,
  left=1mm,
  right=1mm,
  top=1mm,
  bottom=1mm
}

\newtcolorbox{algobox}[1]{
  enhanced,
  breakable,
  colback=green!4,
  colframe=FuseGreen!70,
  title=#1,
  fonttitle=\bfseries,
  boxrule=0.6pt,
  arc=2mm,
  left=1mm,
  right=1mm,
  top=1mm,
  bottom=1mm
}

\newcommand{\code}[1]{\nolinkurl{#1}}
\newcommand{\file}[1]{\nolinkurl{#1}}
\newcommand{\topic}[1]{\nolinkurl{#1}}
\newcommand{\param}[1]{\nolinkurl{#1}}
\newcommand{\csvcol}[1]{\nolinkurl{#1}}
\newcommand{\pkg}[1]{\nolinkurl{#1}}

\title{
  \Huge 多传感器融合算法教程\\
  \Large 面向 Unitree Go2 RSSI 路径进度定位项目
}
\author{Go2 RSSI / Multi-Sensor Fusion Notes}
\date{2026-06-15}

\begin{document}
\maketitle
\tableofcontents

\chapter*{阅读方式}
\addcontentsline{toc}{chapter}{阅读方式}

这份教程的主题是“多传感器融合算法”，但不是泛泛讲 IMU、GPS、雷达融合。它围绕本项目当前 Go2 RSSI 路线复现目标展开：

\begin{center}
  \textbf{用 RSSI 估计机器人在一条预定义路线上的全局路径进度 \(s\)，再用 IMU、odom、\topic{/go2/path_s} 和运动约束让估计更连续、更可信。}
\end{center}

本项目不是二维自由定位优先，而是路径跟踪优先。因此本文把状态空间从一开始就定义在路径坐标上。当前核心状态只有 \(s\)，后续扩展只加入横向偏移 \(e\)：

\[
  \bm{x}_t =
  \begin{bmatrix}
    s_t \\
    e_t
  \end{bmatrix}
\]

其中 \(s_t\) 是沿中心线的路径进度，是当前最重要的输出；\(e_t\) 是相对中心线的横向偏移，是后续扩展。第一版可以只估计 \(s_t\)，数据采集覆盖中心线左右偏移后再扩展到 \((s_t,e_t)\)。

\begin{notebox}{本教程和项目代码的对应关系}
本文重点对照项目中的 baseline 脚本、Phase A 采集说明、真实 \(s\)-only baseline 文档，以及 Go2 RSSI/IMU 采集脚本。尤其是 \file{robot_s_baseline.py}，它已经包含 RSSI 指纹、observed mask、距离度量、confidence、IMU motion gate、滑动滤波、Kalman/EKF 等融合雏形。文中凡是写到滤波状态，均以 \(s\) 或 \((s,e)\) 为准。
\end{notebox}

\chapter{项目中的传感器融合问题}
\label{chap:problem}

\section{真正要融合的是什么}

Go2 RSSI 项目当前有这些信号源：

\begin{footnotesize}
\begin{longtable}{L{0.16\textwidth}L{0.34\textwidth}L{0.40\textwidth}}
\toprule
信号 & 文件/话题 & 在融合中的角色 \\
\midrule
RSSI & \file{rssi_realtime_windows.csv}\newline \file{rssi_raw_stream.csv} & 主要定位观测。它提供“当前位置像哪个路径参考点”，但噪声大、丢包、多径强。 \\
Go2 IMU & \file{imu_stream.csv}，\topic{/go2/imu} & 短时运动状态。适合判断静止/运动、转向、振动、加速度变化，不适合单独长时间积分定位。 \\
Go2 sport state odom & \file{sport_state_straight.csv}，\topic{/odom} & 采集阶段的路径标签来源和评估参考。可用于对齐 RSSI 窗口到 \(s_{gt}\)。 \\
路径进度 & \csvcol{odom_s_m}，\topic{/go2/path_s} & 当前最重要的监督标签和停止条件。它把“机器人走了多远”变成可训练目标。 \\
点云/相机 & \topic{/go2/pointcloud}\newline \file{front_images/} & 当前主要用于安全和复盘，不作为 RSSI 主定位输入。 \\
Nav2 输出 & \topic{/cmd_vel_nav}\newline costmap / planner log & 运动后端/对照信号。Nav2 localization 不应作为 RSSI 主算法输入。 \\
\bottomrule
\end{longtable}
\end{footnotesize}

\section{为什么不是直接估计 \(x,y\)}

室内 RSSI 多径严重，同一个 \(x,y\) 附近的 RSSI 可能随朝向、人体遮挡、机器人姿态变化而变；而路线复现任务天然有一条中心线。用 \(s\) 表示“沿路线走到哪里”有三个优点：

\begin{enumerate}
  \item 标签更容易获得：Go2 运动中可以用 \csvcol{odom_s_m} 或 \topic{/go2/path_s} 对齐 RSSI 窗口。
  \item 控制更直接：路线客户端只需要知道当前 \(s\)，再选择 \(s+\Delta s\) 的局部目标。
  \item 融合更稳定：运动连续性可以自然写成 \(s_t \approx s_{t-1} + \Delta s_t\)。
\end{enumerate}

\section{系统总图}

\begin{figure}[htbp]
\centering
\resizebox{\textwidth}{!}{%
\begin{tikzpicture}[
  node distance=8mm and 12mm,
  every node/.style={font=\small},
  src/.style={draw,rounded corners=1mm,align=center,minimum width=33mm,minimum height=9mm,fill=blue!4},
  sync/.style={draw,rounded corners=1mm,align=center,minimum width=36mm,minimum height=10mm,fill=gray!10},
  alg/.style={draw,rounded corners=1mm,align=center,minimum width=38mm,minimum height=10mm,fill=orange!8},
  outbox/.style={draw,rounded corners=1mm,align=center,minimum width=35mm,minimum height=10mm,fill=green!7},
  arr/.style={-{Latex[length=2mm]},thick},
  faint/.style={-{Latex[length=2mm]},thick,dashed,FuseBlue!70}
]
  \node[src] (rssi) {RSSI window\\mean / std / count};
  \node[src,below=of rssi] (imu) {Go2 IMU\\acc / gyro / rpy};
  \node[src,below=of imu] (odom) {Go2 odom\\\csvcol{odom_s_m} / \topic{/go2/path_s}};

  \node[sync,right=of imu] (syncnode) {时间同步与特征聚合\\RSSI window center time};
  \node[alg,right=of syncnode,yshift=12mm] (meas) {RSSI 测量模型\\KNN / likelihood\\输出 \(z_t^s,c_t\)};
  \node[alg,right=of syncnode,yshift=-12mm] (motion) {运动约束\\IMU gate / odom delta\\输出 \(\Delta s_t,J_t\)};
  \node[alg,right=of meas,yshift=-12mm] (filter) {状态估计\\\(x_t=s_t\) 或 \((s_t,e_t)\)\\EMA / Kalman / EKF};
  \node[outbox,right=of filter] (estimate) {输出\\\(\hat{s}_t,\hat{e}_t,c_t\)};
  \node[outbox,below=of estimate] (control) {路线复现\\direct 或 Nav2 后端};

  \draw[arr] (rssi.east) -- (syncnode.west);
  \draw[arr] (imu.east) -- (syncnode.west);
  \draw[arr] (odom.east) -- (syncnode.west);
  \draw[arr] (syncnode) -- (meas);
  \draw[arr] (syncnode) -- (motion);
  \draw[arr] (meas) -- (filter);
  \draw[arr] (motion) -- (filter);
  \draw[arr] (filter) -- (estimate);
  \draw[arr] (estimate) -- (control);
  \draw[faint] (odom.south east) -- ++(8mm,-5mm) -- ++(82mm,0) -- node[below,font=\scriptsize]{采集/评估标签，不作为 RSSI 主观测} ++(0,16mm);
\end{tikzpicture}%
}
\caption{面向 Go2 路径进度 \(s\) 与后续横向偏移 \(e\) 的融合流程。RSSI 提供绝对但噪声大的路径观测，IMU/odom 提供时间对齐、短时连续性和运动约束。}
\end{figure}

\chapter{数据格式与时间对齐}
\label{chap:data}

\section{RUN\_DIR 输出文件}

每次采集会写到一个 \code{RUN_DIR}。典型文件如下：

\begin{longtable}{p{0.34\textwidth}p{0.56\textwidth}}
\toprule
文件 & 用途 \\
\midrule
\file{rssi_realtime_windows.csv} & RSSI 窗口均值、方差、count。融合算法主要输入。 \\
\file{rssi_raw_stream.csv} & 原始 RSSI 包。用于复盘丢包、尖峰和信标可见性。 \\
\file{imu_stream.csv} & Go2 lowstate IMU：加速度、角速度、roll/pitch/yaw。 \\
\file{sport_state_straight.csv} & Go2 sport state、运动命令、\csvcol{odom_s_m}、点云安全状态。采集阶段最重要的 \(s_{gt}\) 来源。 \\
\file{run_metadata.txt} & 本次采集参数：速度、距离、后端、网卡、RSSI 模式等。 \\
\file{front_images/} & 前置摄像头复盘图像。用于解释异常，不建议第一版加入定位模型。 \\
\bottomrule
\end{longtable}

\section{机器人路径进度标签 \(s\)}

在 Phase A 采集中，\csvcol{odom_s_m} 是从 Go2 sport state 的 \(x,y\) 增量累计得到：

\[
  s_t = \sum_{i=1}^{t} \sqrt{(x_i-x_{i-1})^2 + (y_i-y_{i-1})^2}.
\]

它和命令积分 \csvcol{cmd_s_m} 不一样。\csvcol{cmd_s_m} 只是“我命令机器人应该走了多远”，而 \csvcol{odom_s_m} 更接近“Go2 状态估计认为实际走了多远”。训练 RSSI 指纹库时，优先使用 \csvcol{odom_s_m} 或 ROS 侧 \topic{/go2/path_s} 作为 \(s_{gt}\)。

\section{odom 是什么、怎么来的}

\textbf{odom} 是 odometry（里程计）的缩写，表示机器人在一个局部连续坐标系中的位姿和速度估计。它不是 GPS，也不是全局真实坐标；它通常会在短时间内连续、平滑，但长时间会漂移。

在本项目中，\topic{/odom} 由 Go2 bridge 从 Unitree DDS 的 \code{rt/sportmodestate} 转换而来。Go2 的 sport state 提供机身位置、速度、姿态等估计量；bridge 把这些字段封装为 ROS 2 标准的 \code{nav_msgs/Odometry}，同时发布 \topic{odom -> base_link} TF。这里：

\begin{itemize}
  \item \code{odom} frame 是局部里程计坐标系，起点通常接近本次启动或本次运动的参考位置；
  \item \code{base_link} frame 是机器狗机体坐标系；
  \item \topic{/go2/path_s} 或 \csvcol{odom_s_m} 是从 odom/sport state 位置增量累计出来的路径长度。
\end{itemize}

融合算法中，odom 的角色要分清楚：

\begin{enumerate}
  \item \textbf{采集建库时}：odom/path\_s 可作为 \(s_{gt}\) 标签来源，把 RSSI window 标注到路线进度 \(s\)。
  \item \textbf{在线融合时}：odom 可提供短时运动增量 \(\Delta s^{odom}\) 或连续性参考，但不应替代 RSSI 成为主定位观测。
  \item \textbf{评估复盘时}：odom 可和人工标记、视频、Nav2 日志一起作为误差分析依据。
\end{enumerate}

\begin{warnbox}{odom 不是绝对真值}
如果路线很长、地面打滑、转弯多或 Go2 状态估计漂移明显，odom/path\_s 也会偏。论文或报告中应把它称为采集标签或参考标签，而不是绝对 ground truth。必要时可用人工测距、地面标记、视频或外部定位做校验。
\end{warnbox}

\section{RSSI 窗口与 odom 的时间对齐}

RSSI 与 Go2 状态由不同程序记录，不能假设同一行天然对应。正确做法是按时间对齐：

\begin{lstlisting}
rssi_realtime_windows.csv:
elapsed_sec=12.50, beacon_6_mean=-63

sport_state_straight.csv:
elapsed_sec=12.40, odom_s_m=0.281
elapsed_sec=12.50, odom_s_m=0.284
elapsed_sec=12.60, odom_s_m=0.287

aligned row:
elapsed_sec=12.50, s=0.284, beacon_6_mean=-63
\end{lstlisting}

如果时间戳有 Unix 时间和 steady elapsed 两套，优先用同一台机器、同一采集脚本内部的 \csvcol{elapsed_sec}。跨机器远端采集时，要警惕系统时间不同步；这时可以用脚本启动时间、事件 marker 或运动开始点做校正。

\section{RSSI、IMU、odom 如何同步后进入算法}

实际进入算法的不是“某一瞬间的一条 RSSI”和“某一瞬间的一条 IMU”，而是以 RSSI window 为主时钟组织的一行融合样本。当前算法流程如下：

\begin{enumerate}
  \item RSSI scanner 按窗口输出一行 \file{rssi_realtime_windows.csv}，包括每个 beacon 的 mean/std/count，以及窗口中心时间 \(t_k\)。若文件只有窗口开始时间，可用 \(t_k=t_{start}+0.5\,T_{window}\) 近似。
  \item 在 \file{imu_stream.csv} 中取 \([t_k-\Delta_{imu}/2,\ t_k+\Delta_{imu}/2]\) 内的 IMU 样本，计算 \(\bar{a}_{dyn}\)、\(\bar{\omega}\)、roll/pitch/yaw 波动等派生特征。
  \item 在 \file{sport_state_straight.csv} 或 \topic{/odom} 记录中找离 \(t_k\) 最近的 odom/path\_s，得到采集标签 \(s_{gt}(t_k)\)，也可取前后两点线性插值得到更平滑的标签。
  \item RSSI mean/std/count 进入测量模型，输出 RSSI 对路径进度的观测 \(z_k^s=\hat{s}_{k}^{rssi}\) 和置信度 \(c_k\)。
  \item IMU 派生特征和 odom 增量进入运动约束，输出“这段时间是否真的在动”、允许最大跳变 \(J_k\)、以及先验增量 \(\Delta s_k\)。
  \item 滤波器更新 \(s_k\) 或 \((s_k,e_k)\)，输出当前路径进度估计。
\end{enumerate}

可以把一行融合样本理解为：

\begin{lstlisting}
time t_k:
  RSSI window  -> z_s, confidence
  IMU window   -> motion_score, static/dynamic flag
  odom/path_s  -> label s_gt for training, delta_s prior for online
  filter       -> state estimate s_hat, optional e_hat
\end{lstlisting}

这也是为什么时间戳同步很关键：RSSI window 如果滞后 0.5--1.0 s，而 Go2 以 0.3 m/s 行走，标签就可能错开 0.15--0.30 m；这个误差会直接污染指纹库和在线评估。

\section{最直白版：拿到 IMU 和 RSSI 后怎么得到 \texorpdfstring{\(\hat{s}\)}{s_hat}}

可以把整个融合过程想成四句话。

\begin{enumerate}
  \item \textbf{RSSI 先猜位置。} 当前 RSSI window 会和指纹库里的每个参考点比较。哪个参考点的 RSSI 模式最像当前窗口，算法就先猜当前位置在那个 \(s\) 附近，得到 \(\hat{s}^{rssi}\)。同时算法会给一个 confidence：信标看到得多、最近邻距离小、top-K 候选集中，confidence 就高。
  \item \textbf{IMU/odom 再判断这段时间应该走了多少。} IMU 看加速度和角速度，判断机器狗这 0.5--1 秒是在静止、慢动，还是明显运动。odom 或 \topic{/go2/path_s} 可以给出这段时间的路径增量。于是算法得到一个运动预测：上一帧在 \(s_{t-1}\)，这段时间大概前进 \(\Delta s\)，所以先验位置是 \(s_{prior}=s_{t-1}+\Delta s\)。
  \item \textbf{再决定相信谁多一点。} 如果 RSSI confidence 很高，而且 IMU/odom 也说明机器人确实走了这么远，就让结果更靠近 RSSI 猜测；如果 RSSI 突然跳很远，但 IMU 显示机器人几乎没动，就把结果压回 \(s_{prior}\) 附近。
  \item \textbf{最后输出预测路径进度。} 融合后的结果就是 \(\hat{s}_t\)。它既不是纯 RSSI 的最近邻，也不是纯 odom 积分，而是“RSSI 观测 + IMU/odom 运动预测”共同得到的路径进度。
\end{enumerate}

用最短的伪代码写就是：

\begin{lstlisting}
for each RSSI window:
  s_rssi, confidence = match_rssi_to_fingerprint(rssi_window)
  delta_s = estimate_motion_from_imu_or_odom(imu_window, odom_delta)
  s_prior = previous_s + delta_s
  s_hat = blend(s_prior, s_rssi, confidence, imu_motion_score)
\end{lstlisting}

所以，RSSI 的作用是“告诉算法像哪里”，IMU/odom 的作用是“告诉算法这段时间应该移动多少、RSSI 的跳变合不合理”，最终输出的是融合后的 \(\hat{s}\)。

\section{训练 CSV 的推荐格式}

\file{REAL_S_BASELINE.md} 中定义的真实项目格式是：

\begin{lstlisting}
s,timestamp,beacon_1_mean,beacon_1_std,...,beacon_10_mean,beacon_10_std
\end{lstlisting}

其中：

\begin{itemize}
  \item \csvcol{s} 是全局路线中心线上的真实路径进度，不是每个文件内部单独从 0 开始的局部进度。
  \item 缺失 RSSI 原始记录可用 \code{-999}，算法内部转成弱信号值，例如 \code{-105}。
  \item \csvcol{beacon_i_mean} 是窗口均值；\csvcol{beacon_i_std} 是窗口内波动，可用于权重。
  \item 后续状态扩展只建议加入 \csvcol{e}。若保存 \csvcol{x}/\csvcol{y}，它们应作为复盘/可视化字段，不作为本文主状态。
\end{itemize}

\begin{warnbox}{不要把局部 \(s\) 当成全局 \(s\)}
若某个 CSV 文件里的 \(s\) 是从该文件起点独立累计出来的局部路径长度，就不能把多个文件直接按 \(s\) 排序后当成一条全局路线。Go2 项目必须先定义一条全局中心线，再给每个参考样本统一标注全局 \(s\)。
\end{warnbox}

\chapter{RSSI 指纹定位：融合算法的测量模型}
\label{chap:rssi}

\section{从 Scheme A 开始：纯 RSSI KNN}

最简单的定位模型是：

\[
  z_t^{rssi} \rightarrow \hat{s}_t.
\]

其中 \(z_t^{rssi}\) 是当前窗口的 RSSI 向量：

\[
  \bm{r}_t =
  [r_{1,t}, r_{2,t}, \ldots, r_{M,t}]^\top.
\]

离线建库时，每个参考路径进度 \(s_j\) 对应一个指纹：

\[
  \mathcal{F}_j =
  \{s_j, \bm{\mu}_j, \bm{\sigma}_j, \bm{o}_j\},
\]

其中 \(\bm{\mu}_j\) 是各 beacon 的平均 RSSI，\(\bm{\sigma}_j\) 是标准差，\(\bm{o}_j\) 是 observed rate 或可见性。

\section{加权距离}

\file{robot_s_baseline.py} 中的核心距离形式可以写成：

\[
  d_j(\bm{r}_t)
  =
  \sqrt{
    \frac{
      \sum_{i=1}^{M} w_{j,i}(r_{i,t}-\mu_{j,i})^2
    }{
      \sum_{i=1}^{M} w_{j,i}
    }
  }.
\]

权重来自两个直觉：

\begin{enumerate}
  \item 在某个位置波动小的 beacon 更可信：\(1/\max(\sigma_{j,i},\sigma_{\min})^2\)。
  \item 信号太弱或长期不可见的 beacon 不应该过度影响匹配。
\end{enumerate}

因此代码中使用：

\[
  w_{j,i}
  =
  \text{strength}(\mu_{j,i})
  \cdot
  \frac{1}{\max(\sigma_{j,i}, \sigma_{\min})^2}.
\]

\section{缺失值不是普通数值}

RSSI 缺失很常见。把缺失直接当成 \(-999\) 参与欧氏距离会摧毁匹配，所以项目里采用：

\begin{itemize}
  \item 日志层：缺失可记为 \code{-999}。
  \item 算法层：缺失替换成弱 RSSI，例如 \code{-105}。
  \item mask 层：记录 \csvcol{beacon_i_observed}，区分“真的弱”和“没收到包”。
  \item mismatch 层：训练位置可见、测试不可见，或反过来，要给额外惩罚，但不能等同于普通 RSSI 差值。
\end{itemize}

这对应 \file{robot_s_baseline.py} 的 \code{use_mask}、\code{missing_mismatch_weight}、\code{observed_rate__beacon} 等逻辑。

\section{KNN 输出 \(s\)}

取距离最小的 \(K\) 个状态，用反距离加权：

\[
  \hat{s}_t^{rssi}
  =
  \sum_{j \in \mathcal{N}_K}
  \alpha_j s_j,
  \qquad
  \alpha_j =
  \frac{1/(d_j+\epsilon)}{\sum_{\ell\in\mathcal{N}_K}1/(d_\ell+\epsilon)}.
\]

\begin{algobox}{Scheme A：纯 RSSI 指纹}
\begin{lstlisting}
offline:
  group training rows by global s
  compute beacon mean/std/observed_rate for each s

online:
  read current RSSI window
  compute weighted distance to every fingerprint state
  take top-K nearest states
  output inverse-distance weighted s_pred
\end{lstlisting}
\end{algobox}

\section{距离度量的选择}

\file{robot_s_baseline.py} 已经支持：

\begin{longtable}{p{0.18\textwidth}p{0.32\textwidth}p{0.40\textwidth}}
\toprule
度量 & 形式 & 什么时候考虑 \\
\midrule
Euclidean & 平方误差再开方 & 默认选择，适合 RSSI 差值近似高斯噪声的假设。 \\
Manhattan & 绝对误差 & 对尖峰比欧氏距离更稳健。 \\
Sorensen & 归一化绝对差 & 当不同窗口整体信号强弱漂移明显时可试。 \\
Cosine & 方向相似度 & 更关注 RSSI 模式形状，而不是绝对强度。 \\
\bottomrule
\end{longtable}

第一版实验建议固定 Euclidean 或 Manhattan，不要同时调太多开关。等有稳定 train/test 协议后再系统比较。

\chapter{运动约束：从 Scheme A 到 Scheme B}
\label{chap:motion}

\section{为什么需要运动约束}

RSSI 指纹经常会把两个远处但 RSSI 模式相似的位置混淆。如果每个时刻独立匹配，输出 \(s\) 会跳来跳去。真实机器人运动是连续的，所以应该加入约束：

\[
  s_t \approx s_{t-1} + \Delta s_t.
\]

这就是 Scheme B 的核心。

\section{硬约束：max jump}

如果机器人两次 RSSI window 间隔 \(\Delta t\)，最大速度 \(v_{\max}\)，那么：

\[
  |s_t - s_{t-1}| \le v_{\max}\Delta t + \eta.
\]

代码里的 \param{max_jump_s} 就是这个硬约束。它会把候选状态限制在上一帧附近：

\begin{lstlisting}[language=Python]
constrained = abs(candidate_s - prev_s) <= max_jump_s
\end{lstlisting}

在本项目中，RSSI window 常见参数是 \code{window_sec=1.0}、\code{step_sec=0.5}，Go2 直线采集速度常见 \code{0.30 m/s}。如果每 0.5 s 更新一次，物理上相邻 \(s\) 变化大约 0.15 m。实际可给更宽一点，例如 0.5--1.0 m，用来容忍 RSSI 延迟和标签误差。

\section{软约束：continuity penalty}

硬约束可能过于绝对。软约束是在 RSSI 距离上加一项路径连续代价：

\[
  D_j
  =
  d_j^{rssi}
  +
  \lambda |s_j - (s_{t-1} + \Delta s_{exp})|.
\]

\file{robot_s_baseline.py} 中对应 \param{continuity_lambda} 和 \param{expected_step_s}。

\section{EMA 平滑}

最简单的后处理滤波是指数滑动平均：

\[
  \hat{s}_t
  =
  \alpha \hat{s}_t^{meas}
  +
  (1-\alpha)\hat{s}_{t-1}.
\]

\param{smooth_alpha} 越大，越相信当前 RSSI；越小，越平滑但滞后越明显。

\begin{algobox}{Scheme B：RSSI + 路径连续性 + 平滑}
\begin{lstlisting}
for each RSSI window:
  candidates = all fingerprint states
  if prev_s exists:
      keep states within max_jump_s
      add continuity penalty around prev_s + expected_step_s
  pred_s_raw = KNN(candidates)
  pred_s = alpha * pred_s_raw + (1 - alpha) * prev_s
\end{lstlisting}
\end{algobox}

\section{confidence：什么时候该相信 RSSI}

\file{robot_s_baseline.py} 的 confidence 由四部分组成：

\begin{enumerate}
  \item observed beacon 数量是否足够；
  \item 最近邻距离是否小；
  \item 第一名和第二名距离是否有 margin；
  \item top-K 的 \(s\) 是否集中。
\end{enumerate}

可以写成：

\[
  c_t =
  0.4 c_{obs}
  +0.2 c_{dist}
  +0.2 c_{margin}
  +0.2 c_{spread}.
\]

当 confidence 低时，代码支持把结果往运动先验拉：

\[
  \hat{s}_t
  =
  c_t\hat{s}_t^{rssi}
  +(1-c_t)(\hat{s}_{t-1}+\Delta s_{exp}).
\]

这就是 \param{confidence_blend}。

\chapter{IMU 融合：不是积分位置，而是判断运动状态}
\label{chap:imu}

\section{Go2 IMU 字段}

\file{go2_imu_logger.cpp} 输出：

\begin{lstlisting}
timestamp_iso,timestamp_unix,elapsed_sec,
frame_type,has_acc,has_gyro,has_angle,
acc_x_g,acc_y_g,acc_z_g,
gyro_x_dps,gyro_y_dps,gyro_z_dps,
roll_deg,pitch_deg,yaw_deg
\end{lstlisting}

IMU 最容易让新手踩坑：加速度二次积分得到位置看起来很诱人，但在真实机器人上漂移会很快爆炸。本项目里更适合先把 IMU 用作“短时运动证据”，不是主位置来源。

\section{IMU 派生特征}

\file{robot_s_baseline.py} 中已经实现：

\[
  a_{norm} = \sqrt{a_x^2+a_y^2+a_z^2},
  \qquad
  a_{dyn} = |a_{norm}-1|.
\]

\[
  \omega_{norm}
  =
  \sqrt{\omega_x^2+\omega_y^2+\omega_z^2}.
\]

在 RSSI window 附近取 IMU 均值：

\[
  \bar{a}_{dyn,t}
  =
  \frac{1}{|\mathcal{W}_t|}
  \sum_{\tau \in \mathcal{W}_t} a_{dyn,\tau},
  \qquad
  \bar{\omega}_{t}
  =
  \frac{1}{|\mathcal{W}_t|}
  \sum_{\tau \in \mathcal{W}_t} \omega_{norm,\tau}.
\]

再得到 motion score：

\[
  m_t =
  \text{clip}
  \left(
    \max
    \left(
      \frac{\bar{a}_{dyn,t}}{\theta_a},
      \frac{\bar{\omega}_{t}}{\theta_\omega}
    \right),
    0,1
  \right).
\]

当 \(m_t\) 接近 0，机器人更可能静止或低运动；当 \(m_t\) 接近 1，机器人更可能在明显运动。

\section{IMU motion gate}

代码中 \param{imu_motion_gate} 做两件事。它不是让 IMU 直接输出位置，而是让 IMU 决定“RSSI 的这次跳变是否可信”。

第一，根据 motion score 动态改变最大跳变：

\[
  J_t =
  J_{static}
  +
  m_t (J_{dynamic}-J_{static}).
\]

静止时允许的 \(s\) 跳变小，运动时允许更大。

第二，根据 motion score 调整 RSSI 和上一帧的融合权重：

\[
  \alpha_t =
  \alpha_{static}
  +
  m_t(\alpha_{dynamic}-\alpha_{static}),
\]

\[
  \hat{s}_t =
  \alpha_t \hat{s}_t^{rssi}
  +(1-\alpha_t)\hat{s}_{t-1}.
\]

这非常适合 Go2 场景：RSSI 有时会突然跳到远处相似指纹，但如果 IMU 认为机器人几乎没动，就应该强烈抑制这种跳变。

\section{RSSI 与 IMU 的相互作用}

RSSI 和 IMU 在融合中承担不同角色：

\begin{itemize}
  \item RSSI 提供绝对位置感：当前信号模式像路径上哪个 \(s\) 或 \((s,e)\)。
  \item IMU 提供短时运动证据：这段时间机器人是否移动、是否转身、振动是否异常。
  \item odom/path\_s 提供采集标签和短时先验：训练时标注 \(s_{gt}\)，在线时可作为 \(\Delta s\) 的参考。
\end{itemize}

因此一次更新可以写成三步。

第一步，RSSI 测量模型给出候选：

\[
  z_k^s = h_{rssi}(\bm{r}_k),
  \qquad
  c_k \in [0,1].
\]

其中 \(z_k^s\) 是 KNN/likelihood 输出的路径进度观测，\(c_k\) 是由 beacon 数量、最近邻距离、margin 和 top-K spread 得到的置信度。

第二步，IMU window 给出 motion score：

\[
  m_k = g(\bar{a}_{dyn,k}, \bar{\omega}_k) \in [0,1].
\]

当 \(m_k\) 小，滤波器认为机器人接近静止，RSSI 即使跳到远处也更可能是多径/丢包造成；当 \(m_k\) 大，滤波器允许 \(s\) 有更大变化。

第三步，把 RSSI 观测和运动先验融合：

\[
  \Delta s_k
  =
  \begin{cases}
    \Delta s_k^{odom}, & \text{在线有可靠 odom/path\_s 参考时},\\
    \Delta s_{exp} m_k, & \text{只有 IMU gate 和期望步长时}.
  \end{cases}
\]

\[
  s_{k|k-1} = s_{k-1|k-1} + \Delta s_k.
\]

RSSI 更新时，IMU 和 confidence 共同决定“拉向 RSSI”还是“保持运动先验”：

\[
  \hat{s}_k
  =
  \beta_k z_k^s
  +(1-\beta_k)s_{k|k-1},
  \qquad
  \beta_k = c_k(\alpha_{static}+m_k(\alpha_{dynamic}-\alpha_{static})).
\]

这条公式表达了两个直觉：RSSI 置信度低时少信 RSSI；IMU 判断接近静止时，也少信 RSSI 的远距离跳变。

\begin{algobox}{RSSI + IMU + odom 同步融合}
\begin{lstlisting}
for each RSSI window k centered at t_k:
  # 1. Timestamp alignment
  rssi_features = rssi_window[k]
  imu_features = aggregate_imu(t_k - imu_W/2, t_k + imu_W/2)
  odom_ref = nearest_or_interpolated_odom(t_k)

  # 2. RSSI measurement
  z_s, confidence = knn_or_likelihood(rssi_features)

  # 3. Motion evidence
  motion_score = imu_motion_score(imu_features)
  jump_limit = static_jump + motion_score * (dynamic_jump - static_jump)
  delta_s_prior = odom_delta_or_expected_step(odom_ref, motion_score)

  # 4. State update, state is s or (s, e), not velocity
  s_prior = s_prev + delta_s_prior
  z_s = clamp_to_local_window(z_s, s_prev, jump_limit)
  s_hat = fuse_by_confidence(s_prior, z_s, confidence, motion_score)
\end{lstlisting}
\end{algobox}

\section{IMU 控制输入参数起点}

\begin{longtable}{p{0.30\textwidth}p{0.20\textwidth}p{0.40\textwidth}}
\toprule
参数 & 建议起点 & 含义 \\
\midrule
\param{imu-window-sec} & \code{1.0} & 每个 RSSI window 附近聚合 1 秒 IMU。 \\
\param{imu-static-acc-threshold} & \code{0.04} & 加速度动态分量低于该值偏静止。需按 Go2 数据调整。 \\
\param{imu-static-gyro-threshold} & \code{3.0} & 角速度模长低于该值偏静止。 \\
\param{imu-static-max-jump-s} & \code{0.3--0.5} & 静止/低运动时允许的最大路径跳变。 \\
\param{imu-dynamic-max-jump-s} & \code{1.0--3.0} & 明显运动时允许的最大路径跳变。 \\
\param{imu-static-alpha} & \code{0.10--0.20} & 静止时当前 RSSI 结果权重。 \\
\param{imu-dynamic-alpha} & \code{0.60--0.80} & 运动时当前 RSSI 结果权重。 \\
\bottomrule
\end{longtable}

\chapter{Kalman 与 EKF：把融合写成状态估计}
\label{chap:kalman}

\section{一维路径进度 Kalman Filter}

当只估计路径进度 \(s\) 时，可以定义：

\[
  \bm{x}_t =
  \begin{bmatrix}
    s_t
  \end{bmatrix}.
\]

状态转移：

\[
  \bm{x}_{t|t-1}
  =
  \begin{bmatrix}
    1
  \end{bmatrix}
  \bm{x}_{t-1|t-1}
  +
  \bm{w}_t.
\]

RSSI KNN 输出作为测量：

\[
  z_t = \hat{s}_t^{rssi}
  =
  \begin{bmatrix}
    1
  \end{bmatrix}
  \bm{x}_t
  + n_t.
\]

如果 confidence 低，就增大测量噪声：

\[
  R_t = \frac{R_0}{\max(c_t, 0.05)}.
\]

这已经对应 \file{robot_s_baseline.py} 的 \param{post_filter=kalman}。如果后续扩展到 \((s,e)\)，只需要把状态改成二维：

\[
  \bm{x}_t =
  \begin{bmatrix}
    s_t \\
    e_t
  \end{bmatrix}.
\]

此时 \(e_t\) 表示横向偏移。

\section{Kalman 预测与更新}

预测：

\[
  \hat{\bm{x}}_{t|t-1} = F\hat{\bm{x}}_{t-1|t-1},
  \qquad
  P_{t|t-1} = FP_{t-1|t-1}F^\top + Q.
\]

更新：

\[
  \bm{y}_t = z_t - H\hat{\bm{x}}_{t|t-1},
\]

\[
  S_t = HP_{t|t-1}H^\top + R_t,
  \qquad
  K_t = P_{t|t-1}H^\top S_t^{-1},
\]

\[
  \hat{\bm{x}}_{t|t}
  =
  \hat{\bm{x}}_{t|t-1}
  +
  K_t\bm{y}_t.
\]

\section{EKF：IMU/odom 作为控制输入}

\file{robot_s_baseline.py} 的 \param{post_filter=ekf} 版本更贴近本项目。它先从 IMU/motion score 和 odom 得到一个 \(\Delta s_t^{imu}\) 或 \(\Delta s_t^{odom}\)，再预测：

\[
  u_t = \Delta s_t^{imu/odom},
  \qquad
  s_{t|t-1} = s_{t-1} + u_t,
\]

然后仍用 RSSI KNN 输出的 \(z_t=\hat{s}_t^{rssi}\) 做更新。严格说，当前一维版本更接近带运动先验的 Kalman；称 EKF 是为了给后续加入非线性路线几何、局部曲率或 \((s,e)\) 耦合预留接口。

\begin{algobox}{当前在线融合主循环}
\begin{lstlisting}
for each RSSI window at time t:
  rssi_measurement = KNN(fingerprint_map, rssi_vector)
  confidence = score(observed_dims, distance, margin, topK_spread)
  imu_motion_score = aggregate_imu(t - W/2, t + W/2)
  control_input_u = delta_s_from_imu_or_odom(imu_motion_score, odom_delta)

  predict:
      s_prior = s_prev + control_input_u
      P_prior = F P_prev F^T + Q

  update:
      R = R0 / max(confidence, 0.05)
      s_post = KalmanUpdate(s_prior, rssi_measurement, R)

  publish:
      /path_s_estimate = s_post
\end{lstlisting}
\end{algobox}

\section{参数怎么理解}

\begin{longtable}{p{0.30\textwidth}p{0.20\textwidth}p{0.40\textwidth}}
\toprule
参数 & 代码默认值 & 调参直觉 \\
\midrule
\param{kalman-process-var} & \code{0.05} & 越大越相信运动模型会变化，输出更灵活；越小越平滑。 \\
\param{kalman-measurement-var} & \code{4.0} & RSSI 测量基础噪声。越大越不相信 RSSI，越小越跟随 KNN。 \\
\param{confidence-blend} & off & 开启后低置信度 RSSI 会更靠近运动先验。建议在 RSSI 尖峰明显时测试。 \\
  \param{post-filter} & \code{none} & 可选 \code{sliding_mean}、\code{sliding_median}、\code{kalman}、\code{ekf}。 \\
\bottomrule
\end{longtable}

\chapter{更高级的融合路线}
\label{chap:advanced}

\section{粒子滤波：适合多峰 RSSI}

RSSI 指纹经常多峰：两个不同位置可能 RSSI 很像。Kalman Filter 假设状态近似单峰高斯，而粒子滤波可以保留多个候选。

粒子表示：

\[
  \{s_t^{(n)}, w_t^{(n)}\}_{n=1}^{N}.
\]

预测：

\[
  s_t^{(n)} =
  s_{t-1}^{(n)} + \Delta s_t + \epsilon_t^{(n)}.
\]

观测权重：

\[
  w_t^{(n)}
  \propto
  \exp\left(
    -\frac{d^{rssi}(s_t^{(n)})^2}{2\sigma_r^2}
  \right).
\]

优点是能处理“当前位置可能在 12 m 或 48 m”的模糊；缺点是参数和计算复杂度更高。建议在 Scheme B 和 Kalman 稳定后再做。

\section{HMM/Viterbi：离线轨迹最优}

如果先做离线评估，可以把每个 \(s\) bin 当成一个离散状态，RSSI 距离当观测代价，运动约束当转移代价。Viterbi 求整条轨迹最优：

\[
  \min_{s_{1:T}}
  \sum_t C_{obs}(s_t, z_t)
  +
  \sum_t C_{trans}(s_t, s_{t-1}).
\]

这非常适合分析数据集上“理论上融合能做到多好”，但在线控制时要改成递推版本。

\section{因子图：最终论文级框架}

当系统同时有 RSSI、IMU、odom、路径中心线、可能的 \(e\) 偏移时，因子图会更清晰。这里的优化状态仍限定为 \(\bm{x}_t=[s_t,e_t]^\top\)：

\[
  \min_{\bm{x}_{1:T}}
  \sum_t
  \|r_t - h_{rssi}(\bm{x}_t)\|_{\Sigma_r}^{2}
  +
  \sum_t
  \|\bm{x}_t - f(\bm{x}_{t-1}, u_t)\|_{\Sigma_u}^{2}
  +
  \sum_t
  \|z_t^{odom} - h_{odom}(\bm{x}_t)\|_{\Sigma_o}^{2}.
\]

第一版不需要直接上因子图，但可以把它作为论文叙事的终点：RSSI 是全局但噪声大的观测，IMU/odom 是局部连续约束，路线几何是结构先验。

\section{学习型方法放在哪里}

后续可以用神经网络替换 KNN 的测量模型：

\[
  \hat{s}^{rssi}, c = f_\theta(\bm{r}).
\]

但建议保留同样的融合外壳：网络只输出 RSSI 测量和不确定性，IMU/odom 仍通过滤波器或约束层融合。这样更容易解释，也更容易和 Scheme A/B 公平比较。

\chapter{项目落地流程}
\label{chap:workflow}

\section{采集流程}

推荐逐步采集：

\begin{enumerate}
  \item 定义一条中心线，确定全局 \(s=0\) 到终点。
  \item 固定 beacon 位置；移动 beacon 后必须重新建库。
  \item 先用 direct 后端做 1 m、3 m、10 m 验证。
  \item 正式路线至少正向采 3--5 次。
  \item 用 \csvcol{elapsed_sec} 把 RSSI window 对齐到 \csvcol{odom_s_m}。
  \item 按 \(s\) 分桶或按真实参考点聚合，生成训练 CSV。
  \item 每条测试轨迹单独评估，不要把测试轨迹混成一条伪路径。
\end{enumerate}

\section{运行 baseline}

Scheme A：

\begin{lstlisting}[language=bash]
python3 robot_s_baseline.py \
  --train your_train.csv \
  --test your_test.csv \
  --k 3 \
  --max-jump-s 0 \
  --smooth-alpha 0 \
  --out baseline_outputs/robot_predictions_A.csv \
  --map-out baseline_outputs/robot_map_A.csv
\end{lstlisting}

Scheme B：

\begin{lstlisting}[language=bash]
python3 robot_s_baseline.py \
  --train your_train.csv \
  --test your_test.csv \
  --k 3 \
  --max-jump-s 1.0 \
  --smooth-alpha 0.4 \
  --continuity-lambda 0.2 \
  --expected-step-s 0.15 \
  --out baseline_outputs/robot_predictions_B.csv \
  --map-out baseline_outputs/robot_map_B.csv
\end{lstlisting}

IMU/odom gate + EKF(\(s\))：

\begin{lstlisting}[language=bash]
python3 robot_s_baseline.py \
  --train your_train.csv \
  --test your_test.csv \
  --imu imu_stream.csv \
  --imu-motion-gate \
  --post-filter ekf \
  --confidence-blend \
  --max-jump-s 1.0 \
  --expected-step-s 0.15 \
  --out baseline_outputs/robot_predictions_imu_ekf_s.csv \
  --map-out baseline_outputs/robot_map_imu_ekf_s.csv
\end{lstlisting}

\section{评估指标}

主指标：

\[
  err_t = |\hat{s}_t - s_t^{gt}|.
\]

汇总指标：

\begin{itemize}
  \item mean error；
  \item median error；
  \item p75 / p90 error；
  \item jump count：\(|\hat{s}_t-\hat{s}_{t-1}|\) 超过阈值的次数；
  \item runtime per update；
  \item confidence 与 error 的相关性。
\end{itemize}

论文或报告中第一张对比表建议：

\begin{longtable}{p{0.28\textwidth}p{0.15\textwidth}p{0.15\textwidth}p{0.15\textwidth}p{0.15\textwidth}}
\toprule
方法 & mean & median & p75 & p90 \\
\midrule
RSSI KNN &  &  &  &  \\
RSSI KNN + max jump &  &  &  &  \\
RSSI + continuity + EMA &  &  &  &  \\
RSSI + IMU gate &  &  &  &  \\
RSSI + IMU/odom EKF(\(s\)) &  &  &  &  \\
\bottomrule
\end{longtable}

\section{上线前必须做 shadow mode}

在线闭环前，先让系统只输出 \(\hat{s}\)，但不控制机器狗：

\begin{lstlisting}
RSSI window -> fusion node -> /path_s_estimate
human/video/odom -> compare estimate but do not command Go2
\end{lstlisting}

确认：

\begin{itemize}
  \item \(\hat{s}\) 基本单调或符合运动方向；
  \item 静止时不会大跳；
  \item 丢包时 confidence 下降；
  \item p90 error 在可接受范围内；
  \item 估计异常时运动后端能限速或停止。
\end{itemize}

\chapter{常见问题与调参顺序}
\label{chap:debug}

\section{RSSI-only 误差很大}

优先检查：

\begin{enumerate}
  \item beacon 是否太少或几何分布太差；
  \item 训练和测试时 beacon 是否移动；
  \item 人是否挡在机器人和 beacon 中间；
  \item 缺失值是否正确从 \code{-999} 转成 \param{missing_rssi}；
  \item 是否把多个测试轨迹错误混成一条连续路径；
  \item \(s\) 标签是否是全局路径进度。
\end{enumerate}

\section{\(s\) 估计跳变}

按顺序加：

\begin{enumerate}
  \item \param{max-jump-s}；
  \item \param{continuity-lambda}；
  \item \param{smooth-alpha}；
  \item \param{confidence-blend}；
  \item \param{imu-motion-gate}；
  \item \param{post-filter=kalman} 或 \param{post-filter=ekf}。
\end{enumerate}

\section{输出太平滑，明显滞后}

可能原因：

\begin{itemize}
  \item \param{smooth-alpha} 太小；
  \item \param{kalman-measurement-var} 太大；
  \item \param{max-jump-s} 太小；
  \item IMU motion score 偏低，导致一直按静止处理；
  \item RSSI window 太长，测量本身延迟。
\end{itemize}

\section{IMU gate 没效果}

看 \file{*_predictions.csv} 中这些列：

\begin{itemize}
  \item \csvcol{imu_acc_dynamic_mean}
  \item \csvcol{imu_gyro_norm_mean}
  \item \csvcol{imu_motion_score}
  \item \csvcol{imu_is_static}
  \item \csvcol{effective_max_jump_s}
\end{itemize}

如果 motion score 长期接近 0 或 1，说明阈值不合适。先画出静止段和运动段的 \csvcol{acc_dynamic}、\csvcol{gyro_norm} 分布，再定阈值。

\chapter{附录：文件与参数对照}

\section{关键代码函数}

\begin{longtable}{p{0.32\textwidth}p{0.58\textwidth}}
\toprule
函数 & 作用 \\
\midrule
\code{load_csvs()} & 读取训练/测试 CSV，处理 \code{-999} 缺失，生成 observed mask。 \\
\code{build_fingerprint_map()} & 按全局 \(s\) 聚合 beacon mean/std/observed rate，生成指纹库。 \\
\code{make_weights()} & 根据 RSSI 强度和 std 计算 beacon 权重。 \\
\code{compute_distances()} & 支持 Euclidean、Manhattan、Sorensen、Cosine 距离。 \\
\code{compute_confidence()} & 根据可见 beacon、距离、margin、top-K spread 计算置信度。 \\
\code{attach_imu_features()} & 把 IMU window 聚合到每个 RSSI row，生成 motion score。 \\
\code{predict_s()} & 主融合循环：KNN、连续性、confidence、IMU gate、后处理滤波。 \\
\code{summarize_errors()} & 输出 mean/median/p75/p90。 \\
\bottomrule
\end{longtable}

\section{建议阅读的项目资料}

\begin{itemize}
  \item \file{REAL_S_BASELINE.md}
  \item \file{PROJECT_PLAN_AB.md}
  \item \file{RSSI_BASELINE.md}
  \item \file{GO2_PHASE_A_STRAIGHT_LINE.md}
  \item \file{GO2_NAV2_INTEGRATION.md}
  \item \file{robot_s_baseline.py}
  \item \file{rssi_fingerprint_baseline.py}
  \item \file{scripts/go2_collect_rssi_imu.sh}
  \item \file{go2_sdk/go2_imu_logger.cpp}
\end{itemize}

\end{document}
